\cite{briou2021trends} talks lots about different ways to do DA cycloadditions. talks about electronic effects (2.3)


\cite{ratwani2023self} This is a recent reivew talking about different DA for use in soft robotics ( states that cp is nice because repeating unit but not very strong)


\cite{kuhl2016self} michael addition sulfur stuff

\cite{terryn2021review} the paper with that comparison table


\cite{mei2016highly} pi stacking with pt -pt interactiosn


\cite{yang2018self} supramolecular stuff (review

\cite{burattini2009novel} specific example of pi stacked self healing


\cite{srivastav2025comprehensive} review about self healin polyerms


\cite{guo2025diels}  2025 paper about diels alder polymers

\cite{kashef2025toward} 2025 review of soft robots SH

\cite{cioffi2022review} a review about using sH polymer for struturla things


\cite{terryn2017self} SH soft pneu robotics review

%\cite{26329} deterioration of polyimide through water decomposition

\cite{moore1964gel} GPC original paper


\cite{izunobi2011polymer} NMR polymer determination

\cite{shit1986application}NMR for MWD


Advce on formatting a report https://chemlab.truman.edu/the-laboratory-report/#Experimental


https://doi.org/10.1002/pol.1957.1202310436 : amide hydrolysis in polyvinylamide. The alkaline hydrolysis is second order (this has citations) 

DOI 10.1002/poc.1492 :  information on secondary amide hydrolysis. this has plenty of references about the 

\cite{o1970acidic} : Big paper with stuff about amide hydrolysis. 



% \cite{example}
% Purpose / Research Question(s): Example: To examine how AI-powered feedback affects student performance in writing.
% Theoretical Framework (if applicable): Example: Constructivist learning theory.
% Methodology: Example: Quantitative, experimental design.
% Participants & Context: Example: 120 high school students (ages 14-17), across two schools.
% Data Collection Tools: Example: Pre- and post-test, Likert-scale survey.
% Data Analysis Method: Example: T-tests for writing score comparison.
% Key Findings: Example: Students in the experimental group showed a 15% improvement in writing scores.
% Notable Quotes: Example: “AI feedback provides students with immediate, personalized guidance.”
% Limitations / Biases: Example: Small sample size; potential bias in self-reported feedback.


\cite{dow2010modification}
Purpose: Kapton hydrolysis
MEthods: Reaction of KOH with PI to make PAA
Reagents: 5M KOH solution at 50C (varied times)
Chgaracterisation: IR
IR of Kapton:  Several representative absorbances were observed at 1775 cm1 (imide C]O, asymmetric stretch), 1720 cm1 (imide C]O, symmetric stretch), 1500 cm1 (benzene ring, C]Cstretch), 1370 cm1 (imide C–N–C, ring stretch), and 1235 cm1 (aromatic ether, Ar–O–Ar stretch)
After hydrolyis IR: After the hydrolysis reaction with KOH, the absorbance peaks at 1775 cm1, 1720 cm1, and 1370 cm1 disappeared, but the bands at 1500 cm1 and 1235 cm1 remained. On the other hand, several new absorbance peaks appeared at 1640 cm1 (amide, N–H stretch), 1600 cm1 (carboxylate salt, asymmetric stretch), 1545 cm1 (amide, N–H bend), 1410 cm1 (carboxylate salt, symmetric stretch), and 1220 cm1 (carboxylic acid, C–O)

\cite{liu2019efficient}
Purpose: Kapton hydrolysis
MEthods: Activation through grafting of COOH
IR Characterisation of Kapton: "onsists of a peak at 1370 cm−1 for C−N stretching and out-of-plane bending of C−N−C and two peaks at 1720 and 1780 cm−1 for CO symmetrical and asymmetrical stretching of the imide ring."
They include an IR plot with before and after KOH hydrolysis where the difference between the lines isnt very significant (2000 and lower) "three characteristic peaks at 1720, 1500, and 1370 cm−1can be used to determine the extent of the imide ring cleavage" and "the peak at 1720 cm−1represents the symmetric stretching C=O bond of the imide ring, and its intensity reduces as the imide ring cleaves under the KOH attack."

\cite{liu2019efficient}
Pupose: kapton casting
Methods: dried (solvent = dmac) on a glass sheet at 80C for 12hr and then 100, 200 and 300C for 1 hour each


\cite{wu2017study}
Purpose: KApton hydrolyis
Methods: KOH (0.1,0.5,1,5 M solutions) 1-15mins at room temp
Characterisation: Compares the ratio of the 1720:1500 peaks in the Kapton IR before and after hydrolysis. They found that this ratio decreased with increasing KOH concentration, ( they used this to estimate the number of imide rings that had been broken)
Key obs: below 0.5M KOH this was pretty inneffective


\cite{chang2000two}
Purpose: Kapton hydrolysis

\cite{chang2000two}
Purpose: Kapton casting
Methods: cast on a glass plate with solvent(DMAc) evaporated in vacuo at 50C for 6hrs, this was dried again for 50C 6hrs out of vacuum and then samples of PAA were heated under vacuum for 30 min at 300, 400, 420, 430, 440, 450, and 550 °C (above 400-450C a rearrangement happens to form something else )
Analysis: [WARNING THIS IS NOT STRICTLY KAPTON THE AR RINGS HAVE OH GROUPS ON] before heating PAA exhibited broad IR bands at 2500-3700(O-H and N-H stretches) and also 1650 (amide c=o stretch). After heating at 300, the ap-pearance of characteristic imide carbonyl peaksat 1783 (symmetric C=O stretch), 1717 (asym-metric C=O stretching), and 720 cm-1 (imide ring bend.)
Observations: PAA was readily sol-uble in a strong acid and in polar aprotic solvents such as DMAc, NMP, and DMSO at room temper-ature. However, none of the heat-treated poly-mers were soluble at all


\cite{xue1986studies}
Purpos: hydrolysis
methods: samples immersed into 0.25N ( N for normality) KOH solution at 30C for 55seconds, substrate removed and acidified with HCl. THere was a gel remaining of the surface of the kapton which was dissolved in DMF and then removed, acidified and then when water was added a precipitate formed and then IR was done.
Analysis: IR The band at 1780 cm-' is due to the imide residue. The bands at 1660 and 1540 cm-' are due to the amide structure. The band at 1722 cm-' is assigned to the carbonyl stretching mode of the carboxylic acid. The band at 1500 cm-' is due to the phenyl mode which can be used as an internal thickness band. [they also did a comparison of the 1720 and 1500 ratios]
Observation: they found that adding a divalent salt (something with a 2+ (zncl2 etc)) slowed the rate of hydrolysis.

\cite{thomas1992organic}
Purpose: hydrolysis
Methods: samples prepped from PAA by spin coating a a layer of PAA onto a single crystal silicon substrate with a sputtered layer of chromium foradhesion. They were dried at 80C for 30 mins. The converted to kapton by heating at 150C 30 min, 200C 30min, 300C 30 min and 400C 60 mion in a vacuum oven. Hydrolysis used 0.25M NaOH -> based on results they think that hydrolysis occurs at 160angstroms/hour. THey calculated this by looking at the c=O imide strecthing peak at 1740. "Sincethereare anumberofpeakspresent intheregion1500 1740cm-1, including1727cm-1(C=0stretchof theCO2Hgroup), 1670cm-1 (C=0stretchof theCONHgroup),1606cm-1(C=0stretchof theCO2-groupandthephenylring(of thePMDAmoiety)stretch),1583cm-1(NHbendof theCONHgroup),1501and1513cm-1(crystallineandamorphousphenyl ring(of theODAmoiety) stretches)"

\cite{stoffel1996surface}
Purpose: hydrolysis
MEthods: uses a range of bases (1M solutions of KOH,CsOH,RbOH,and(CH3)4NOH.) and then acetic acid to remove the metal ion and convert into PAA. Peuieces of film were cut to approx 1 inch square and dipped into the base for measured time. Metal ions are then detectable by rtherford backscattering, with (CH3)4NOH, no ions were detected
Observations: At room temperature the rate of hydrolysis with KOH is rather slow and this paper has full data to back it up. THey studied between 25 and 60 C and upto a time of 100mins. They found that the activation energy with KOH was 60KJ/mol and woth (CH3)4NOH it was 67kj/mol

\cite{plechaty1992chemical}
Purpoise: hydrolysis
MEthods: Using 0.25M NaOH for 2-5hrs at room temp, then treatment with 0.1M acetic acid
analysis: with xray photoelectron spectroscopy

\cite{pawlowski1991etch}
